\section{Setup}
\label{sec:setup}

% In this section, we present the list of text-to-video generative models benchmarked on the \name{} dataset and provide further details about the dataset splits. 

\paragraph{Video Generative Models.}
\label{sec:t2v_generative_model}

We evaluate a diverse range of \textbf{twelve} closed and open text-to-video generative models on \name{} dataset. The list of the models includes \textit{ZeroScope} \cite{zeroscope}, \textit{LaVIE} \cite{wang2023lavie}, \textit{VideoCrafter2} \cite{chen2024videocrafter2}, \textit{OpenSora} \cite{OpenSora}, CogVideoX-2B and 5B \cite{yang2024cogvideox}, \textit{StableVideoDiffusion (SVD)-T2I2V} \cite{blattmann2023stable}, \textit{Gen-2 (Runway)} \cite{gen2}, \textit{Lumiere-T2V}, \textit{Lumiere-T2I2V} (Google) \cite{bar2024lumiere}, Dream Machine (Luma AI) \cite{lumalabsLumaDream}, and \textit{Pika} \cite{pika}. 
% Here, the \textit{T2I2V} models involve the generation of an image (\textit{I}) conditioned on the caption (\textit{T}) followed by video generation (\textit{V}) conditioned on the generated image. 
We provide more model and inference details in Appendix \ref{app:models} and \ref{app:inference_details}. \footnote{While there are various closed models such as Sora \cite{liu2024sora}, Kling AI \cite{klingai}, and Genmo \cite{genmo}, we could not get access through their videos due to the lack of API support.}


\paragraph{Dataset setup.}
\label{sec:setup_dataset}


As described earlier, we train \model{} to enable cheaper and scalable testing of the generated videos on our dataset (§ \ref{sec:automatic_evaluation}). To facilitate this, we split the prompts in the \name{} dataset equally into \textit{train} and \textit{test} sets. Specifically, we utilize the human annotations on the generated videos for the $344$ prompts in the \textit{test} set for benchmarking, while the human annotations on the generated videos for the $344$ prompts in the \textit{train} set are used for training the automatic evaluation model. We ensure that the distribution of the state of matter (solid-solid, solid-fluid, fluid-fluid) and complexity (easy, hard) is similar in the training and testing.

\paragraph{Benchmarking.} Here, we generate one video per test prompt for each T2V generative model in our testbed. Subsequently, we ask three human annotators to judge the semantic adherence and physical commonsense of the generated videos. In our experiments, we report the majority-voted scores from the human annotators. We find that the inter-annotator agreement for semantic adherence and physical commonsense judgment is $75\%$ and $70\%$, respectively. This indicates that the human annotators find the task of judging physical commonsense more subjective than semantic adherence. \footnote{Variations in annotations arise from differing tolerance for commonsense violations in imperfect videos. As generative models improve, human annotations will align more closely.} In total, we collect $\hb{24500}$ human annotations across the testing prompts and T2V models.


\paragraph{Training set for \model{}.} Here, we sample two videos per training prompt for nine T2V models.\footnote{\hb{Since the CogVideoX and Dream Machine models were released very recently, they could not be included in the training set of the automatic evaluator.}} We choose two videos to obtain more data instances for training the automatic evaluation model. Subsequently, we ask one human annotator to judge the semantic adherence and physical commonsense of the generated videos. In total, we collect $12000$ human annotations, half of them for semantic adherence and the other half for physical commonsense. Specifically, we finetune \videocon{} to maximize the log likelihood of \textit{Yes}/\textit{No} conditioned on the multimodal template for semantic adherence and physical commonsense tasks (Appendix \ref{app:videocon_details}). We do not collect three annotations per video as it is financially expensive. In total, we spent $\$\hb{3500}$ on collecting human annotations for benchmarking and training.

